% ============================================================
\chapter{Get Code from Online Repositories. How to Use GitHub}
\label{ch:cli-scripting}
% ============================================================

% ============================================================
\section{Configuring SSH Authentication for GitHub on Linux}
% ============================================================
% ------------------------------------------------------------
\subsection{Configuring SSH Authentication for GitHub on Linux}
% ------------------------------------------------------------
GitHub no longer supports password authentication for Git operations over HTTPS. 
Instead, authentication must be performed using either a Personal Access Token or 
an SSH key. The recommended method for Linux environments is SSH authentication.

The following procedure generates an SSH key and configures Git to authenticate
with GitHub.

\subsubsection{Generate an SSH Key}

Create a new SSH key using the \texttt{ed25519} algorithm:

\begin{lstlisting}[basicstyle=\footnotesize\ttfamily]
ssh-keygen -t ed25519 -C "your_email@example.com"
\end{lstlisting}

When prompted for the file location, press \texttt{Enter} to accept the default:

\begin{lstlisting}[basicstyle=\footnotesize\ttfamily]
~/.ssh/id_ed25519
\end{lstlisting}

Optionally provide a passphrase for additional security.

\subsubsection{Start the SSH Agent}

Start the SSH agent and add the newly created key:

\begin{lstlisting}[basicstyle=\footnotesize\ttfamily]
eval "$(ssh-agent -s)"
ssh-add ~/.ssh/id_ed25519
\end{lstlisting}

\subsubsection{Add the Public Key to GitHub}

Display the public key:

\begin{lstlisting}[basicstyle=\footnotesize\ttfamily]
cat ~/.ssh/id_ed25519.pub
\end{lstlisting}

Copy the entire output and add it to GitHub:

\begin{enumerate}
\item Navigate to \texttt{GitHub → Settings → SSH and GPG keys}
\item Click \texttt{New SSH key}
\item Paste the copied key and save
\end{enumerate}

\subsubsection{Verify the Connection}

Test the authentication:

\begin{lstlisting}[basicstyle=\footnotesize\ttfamily]
ssh -T git@github.com
\end{lstlisting}

A successful configuration will display a message similar to:

\begin{lstlisting}[basicstyle=\footnotesize\ttfamily]
Hi <username>! You've successfully authenticated.
\end{lstlisting}

\subsection{Clone a Repository Using SSH}

Repositories must be cloned using the SSH URL:

\begin{lstlisting}[basicstyle=\footnotesize\ttfamily]
git clone git@github.com:USER/REPOSITORY.git
\end{lstlisting}

Using SSH avoids credential prompts and enables secure automated
Git operations.

% ============================================================
\subsection{Pushing to Multiple Remote Repositories}
\label{subsec:multiple-remotes}
% ============================================================

A single local repository can maintain simultaneous connections to any
number of remote repositories. This is useful when a project must be
mirrored across an organisation account and a personal account---for
example, keeping the \textsc{Dair-3} workshop materials synchronised
between the \texttt{DAIR3} organisation and the personal
\texttt{biomathematicus} GitHub account.

\subsubsection{Concepts}

Git identifies remote repositories by short \emph{names}. The name
\texttt{origin} is the conventional default assigned when a repository
is first cloned. Additional remotes can be registered under any name.
Each remote entry stores a fetch URL and a push URL, both of which can
be set independently.

\begin{infobox}
A remote name is purely local---it exists only in the
\texttt{.git/config} file on your machine and is never pushed to any
server. Two developers cloning the same repository can give the same
remote entirely different names without consequence.
\end{infobox}

\subsubsection{Listing Registered Remotes}

Before adding a remote, confirm what is already registered:

\begin{lstlisting}[style=bashstyle, caption={List all remotes and their URLs}]
git remote -v
\end{lstlisting}

Each remote appears twice---once for fetch and once for push. A freshly
cloned repository typically shows only \texttt{origin}:

\begin{lstlisting}[style=inlineStyle]
origin  https://github.com/DAIR3/DAIR3-Workshop.git (fetch)
origin  https://github.com/DAIR3/DAIR3-Workshop.git (push)
\end{lstlisting}

\subsubsection{Adding a Second Remote}

Register an additional remote with \texttt{git remote add}:

\begin{lstlisting}[style=bashstyle, caption={Register a personal mirror remote (Linux / WSL)}]
git remote add biomathematicus \
    https://github.com/biomathematicus/DAIR3-Workshop.git
\end{lstlisting}

\begin{lstlisting}[style=psstyle, caption={Register a personal mirror remote (PowerShell)}]
git remote add biomathematicus `
    https://github.com/biomathematicus/DAIR3-Workshop.git
\end{lstlisting}

Verify the result:

\begin{lstlisting}[style=bashstyle]
git remote -v
\end{lstlisting}

The output should now show four lines:

\begin{lstlisting}[style=inlineStyle]
biomathematicus  https://github.com/biomathematicus/DAIR3-Workshop.git (fetch)
biomathematicus  https://github.com/biomathematicus/DAIR3-Workshop.git (push)
origin           https://github.com/DAIR3/DAIR3-Workshop.git (fetch)
origin           https://github.com/DAIR3/DAIR3-Workshop.git (push)
\end{lstlisting}

\subsubsection{Pushing to Each Remote}

With two remotes registered, every push must be issued explicitly to
each destination. Git does not push to multiple remotes automatically
unless a push target is configured:

\begin{lstlisting}[style=bashstyle, caption={Push to both remotes explicitly}]
git push origin HEAD
git push biomathematicus HEAD
\end{lstlisting}

\texttt{HEAD} resolves to the current branch, which is the safest form
to use in scripts because it does not depend on knowing the branch name
in advance.

\subsubsection{Checking Whether a Remote Exists Before Pushing}

In automated scripts it is good practice to verify that a remote is
registered before attempting to push to it, so a missing remote produces
a clear diagnostic rather than a fatal error.

\begin{lstlisting}[style=bashstyle, caption={Guard a push with a remote existence check (bash)}]
if git remote get-url biomathematicus > /dev/null 2>&1; then
    git push biomathematicus HEAD
else
    echo "Remote 'biomathematicus' not configured."
    echo "Register it with:"
    echo "  git remote add biomathematicus \
https://github.com/biomathematicus/DAIR3-Workshop.git"
fi
\end{lstlisting}

\begin{lstlisting}[style=psstyle, caption={Guard a push with a remote existence check (PowerShell / batch)}]
git remote get-url biomathematicus >nul 2>&1
if errorlevel 1 (
    echo Remote 'biomathematicus' not configured.
    echo Register it with:
    echo   git remote add biomathematicus ^
        https://github.com/biomathematicus/DAIR3-Workshop.git
) else (
    git push biomathematicus HEAD
)
\end{lstlisting}

\subsubsection{Managing Remotes}

Table~\ref{tab:git-remote-commands} summarises the full set of remote
management commands.

\begin{table}[h]
\centering
\caption{Git remote management commands}
\label{tab:git-remote-commands}
\begin{tabular}{lp{8.5cm}}
\toprule
Command & Effect \\
\midrule
\texttt{git remote -v} &
  List all remotes with their fetch and push URLs. \\
\texttt{git remote add <name> <url>} &
  Register a new remote under the given name. \\
\texttt{git remote remove <name>} &
  Delete a remote entry from the local configuration. \\
\texttt{git remote rename <old> <new>} &
  Rename a remote without changing its URL. \\
\texttt{git remote set-url <name> <url>} &
  Update the URL of an existing remote. \\
\texttt{git remote get-url <name>} &
  Print the URL of a named remote; exits with an error if
  the remote does not exist (useful in scripts). \\
\texttt{git remote show <name>} &
  Display detailed information about a remote, including
  tracked branches and push configuration. \\
\bottomrule
\end{tabular}
\end{table}

\subsubsection{Alternative: a Single Remote with Multiple Push URLs}

Git also supports configuring one remote name to push to two URLs
simultaneously, so that a single \texttt{git push origin HEAD} reaches
both destinations. This is convenient but reduces visibility---a failure
on the second URL can be easy to miss in the output.

\begin{lstlisting}[style=bashstyle, caption={Add a second push URL to the origin remote}]
# The first push URL is already set from the clone.
# Add the personal repo as an additional push target:
git remote set-url --add --push origin \
    https://github.com/DAIR3/DAIR3-Workshop.git
git remote set-url --add --push origin \
    https://github.com/biomathematicus/DAIR3-Workshop.git
\end{lstlisting}

\begin{warningbox}
When using \texttt{set-url --add --push}, note that the \emph{first}
call replaces the default push URL rather than adding to it. Both URLs
must be specified explicitly as shown above. Verify the result with
\texttt{git remote -v} before relying on this configuration in
automated scripts.
\end{warningbox}

For the \textsc{Alice} project, the explicit two-remote approach
(separate \texttt{git push} calls per remote) is preferred because
\texttt{gh.bat} checks \texttt{errorlevel} after each push and can
report failures independently. A combined push URL would surface only
the second failure.

% ============================================================
\section{Python Virtual Environments}
\label{sec:venvs}
% ============================================================

Before we write any scripts, we need to understand \emph{what} those scripts
will manage: Python virtual environments.

\textbf{The Problem: Package Conflicts}. 
Imagine you have two projects on your computer. Project~A uses version~1.5
of a library called \texttt{numpy}; Project~B requires version~2.0. If you
install both versions globally, they will conflict. The last one installed
wins, and one of your projects will break.

\textbf{The Solution: Virtual Environments}.  
A \textbf{virtual environment} (or \emph{venv}) is an isolated copy of
Python with its own set of installed packages, completely separate from
every other environment on your machine. Think of it as a self-contained
bubble around your project.

\begin{itemize}
  \item Each project gets its own environment
  \item Packages installed in one environment cannot interfere with another
  \item You can recreate an environment exactly from a list of packages
        (\texttt{requirements.txt})
  \item Deleting an environment is safe and clean
\end{itemize}

\subsection{Creating a Virtual Environment}

\subsubsection{Windows (CMD or PowerShell)}

\begin{lstlisting}[style=windowsStyle, caption={Creating a venv on Windows}]
:: Navigate to your project folder
cd C:\Projects\DAIR3

:: Create the environment (Python must be installed globally)
python -m venv C:\penv\DAIR3
\end{lstlisting}

\subsubsection{Linux / WSL}

\begin{lstlisting}[style=linuxStyle, caption={Creating a venv on Linux/WSL}]
# Navigate to your project folder
cd ~/DAIR3

# Create the environment
python3 -m venv ~/penv/DAIR3
\end{lstlisting}

In both cases, \texttt{C:\textbackslash penv\textbackslash DAIR3} or
\texttt{\textasciitilde/penv/DAIR3} is where the environment will be stored.
Keeping all environments in a single folder (\texttt{penv} stands for
``Python environments'') makes them easy to find and manage.

\subsection{Activating a Virtual Environment}

Activating an environment modifies your current shell session so that
\texttt{python} and \texttt{pip} refer to the isolated copies inside the
venv, not the global ones.

\subsubsection{Windows}

\begin{lstlisting}[style=windowsStyle, caption={Activating a venv on Windows}]
C:\penv\DAIR3\Scripts\activate.bat
\end{lstlisting}

After activation, your prompt will change to show the environment name:
\begin{lstlisting}[style=windowsStyle]
(DAIR3) C:\Projects\DAIR3>
\end{lstlisting}

\subsubsection{Linux / WSL}

\begin{lstlisting}[style=linuxStyle, caption={Activating a venv on Linux/WSL}]
source ~/penv/DAIR3/bin/activate
\end{lstlisting}

After activation:
\begin{lstlisting}[style=linuxStyle]
(DAIR3) drg@machine:~/DAIR3$
\end{lstlisting}

\begin{warningbox}
Note the \texttt{source} keyword in the Linux command. This is critical and
is explained in detail in Section~\ref{sec:source-vs-execute}.
\end{warningbox}

\subsection{Installing Packages}

Once the environment is active, use \texttt{pip} to install packages.
The syntax is identical on both platforms:

\begin{lstlisting}[style=inlineStyle, caption={Installing packages with pip}]
pip install numpy pandas matplotlib openai langchain
\end{lstlisting}

To save a list of all installed packages (so you or a colleague can
recreate the environment later):

\begin{lstlisting}[style=inlineStyle]
pip freeze > requirements.txt
\end{lstlisting}

To recreate an environment from such a list:

\begin{lstlisting}[style=inlineStyle]
pip install -r requirements.txt
\end{lstlisting}

\subsection{Deactivating}

To leave the environment and return to the global shell (identical on
both platforms):

\begin{lstlisting}[style=inlineStyle]
deactivate
\end{lstlisting}

% ============================================================
\section{Level 1 --- Script Files}
\label{sec:level1}
% ============================================================

Now that we understand virtual environments, we can write scripts to
activate them automatically. The goal is to create two scripts:

\begin{enumerate}
  \item \textbf{\texttt{ActivateEnv}} --- A reusable script that accepts
        two parameters (environment name and project directory), activates
        the environment, and opens VS Code.
  \item \textbf{\texttt{DAIR3}} --- A project-specific one-liner that calls
        \texttt{ActivateEnv} with the correct parameters for the DAIR3
        project.
\end{enumerate}

\subsection{The Windows Solution: Batch Files}

\subsubsection{ActivateEnv.bat}

The Windows version, \texttt{ActivateEnv.bat}, is stored in
\texttt{C:\textbackslash penv} and placed on the system \texttt{PATH}.
It accepts two parameters: \texttt{\%1} (environment name) and
\texttt{\%2} (project directory).

\begin{lstlisting}[style=windowsStyle, basicstyle=\footnotesize\ttfamily, caption={\texttt{ActivateEnv.bat} --- Reusable environment launcher (Windows)}]
@echo off
setlocal

:: Read parameters
set "ENV_NAME=%~1"
set "BASE_DIR=%~2"
echo Starting activation of environment '\%ENV_NAME%'...

:: Validate input
if "%ENV_NAME%"=="" (
    echo [ERROR] Environment name not specified.
    goto :eof
)
if "%BASE_DIR%"=="" (
    echo [ERROR] Base directory not specified.
    goto :eof
)

:: Check if base directory exists
if not exist "%BASE_DIR%" (
    echo [ERROR] Directory not found: %BASE_DIR%
    goto :eof
)

:: Change to working directory
cd /d "%BASE_DIR%"

:: Define activation script path
set "VENV_ROOT=C:\penv"
set "VENV_ACTIVATE=%VENV_ROOT%\%ENV_NAME%\Scripts\activate.bat"

:: Activate the environment
if exist "%VENV_ACTIVATE%" (
    echo Activating...
    call "%VENV_ACTIVATE%"
    echo Activation complete.
) else (
    echo [ERROR] Not found: %VENV_ACTIVATE%
    goto :eof
)

:: Open VS Code in the project directory
echo Launching VS Code...
code .

endlocal
\end{lstlisting}


Key Windows-specific elements:
\begin{itemize}
  \item \texttt{@echo off} suppresses the echoing of each command to the
        screen.
  \item \texttt{setlocal} ensures that variable changes are local to this
        script.
  \item \texttt{\%\textasciitilde1} strips surrounding quotes from the first
        argument.
  \item \texttt{call} is used to invoke another batch file and return
        control afterward.
  \item \texttt{goto :eof} jumps to the end of the file (a clean exit).
\end{itemize}

\subsubsection{DAIR3.bat}

This file lives in the DAIR3 project directory. It simply calls
\texttt{ActivateEnv.bat} with the correct arguments:

\begin{lstlisting}[style=windowsStyle, caption={\texttt{DAIR3.bat} --- Project launcher (Windows)}]
ActivateEnv DAIR3 "C:\Projects\DAIR3"
\end{lstlisting}

Because \texttt{C:\textbackslash penv} is on the system \texttt{PATH},
Windows finds \texttt{ActivateEnv.bat} automatically. You can run
\texttt{DAIR3.bat} by double-clicking it or typing \texttt{DAIR3} in any
CMD window.

\subsection{The Linux Solution: Shell Scripts}
\label{subsec:linux-scripts}

\subsubsection{Understanding Shell Script Syntax}

A Linux shell script is a plain text file containing Bash commands. Every
script should begin with a \textbf{shebang line}---a special first line that
tells the system which interpreter to use:

\begin{lstlisting}[style=linuxStyle]
#!/bin/bash
\end{lstlisting}

Parameters passed to a script are accessed as \texttt{\$1}, \texttt{\$2},
\texttt{\$3}, etc. The script's own name is \texttt{\$0}.

\subsubsection{ActivateEnv.sh}

\begin{lstlisting}[style=linuxStyle, basicstyle=\footnotesize\ttfamily, caption={\texttt{ActivateEnv.sh} --- Reusable environment launcher (Linux/WSL)}]
#!/bin/bash
# ActivateEnv.sh
# Usage: source ~/penv/ActivateEnv.sh <ENV_NAME> <PROJECT_DIR>
# Note: Must be *sourced*, not executed. See explanation below.

ENV_NAME="$1"
BASE_DIR="$2"

echo "Starting activation of environment '${ENV_NAME}'..."

# Validate input
if [ -z "$ENV_NAME" ]; then
    echo "[ERROR] Environment name not specified."
    return 1
fi

if [ -z "$BASE_DIR" ]; then
    echo "[ERROR] Base directory not specified."
    return 1
fi

# Check if the project directory exists
if [ ! -d "$BASE_DIR" ]; then
    echo "[ERROR] Directory not found: ${BASE_DIR}"
    return 1
fi

# Change to the project directory
cd "$BASE_DIR" || return 1

# Build the path to the activation script
VENV_ROOT="$HOME/penv"
VENV_ACTIVATE="${VENV_ROOT}/${ENV_NAME}/bin/activate"

# Activate the environment
if [ -f "$VENV_ACTIVATE" ]; then
    echo "Activating..."
    # shellcheck source=/dev/null
    source "$VENV_ACTIVATE"
    echo "Activation complete."
else
    echo "[ERROR] Not found: ${VENV_ACTIVATE}"
    return 1
fi

# Open VS Code in the project directory (WSL-aware)
echo "Launching VS Code..."
code .
\end{lstlisting}

Notice several differences from the Windows version:
\begin{itemize}
  \item \texttt{-z} tests whether a string is empty (zero length).
  \item \texttt{-d} tests whether a path is a directory.
  \item \texttt{-f} tests whether a path is a regular file.
  \item \texttt{\$HOME} is the Linux equivalent of
        \texttt{\%USERPROFILE\%} --- it expands to your home directory.
  \item \texttt{return 1} exits the script with an error code (we use
        \texttt{return} rather than \texttt{exit} for a reason explained in
        the next section).
\end{itemize}

\subsubsection{DAIR3.sh}

\begin{lstlisting}[style=linuxStyle, caption={\texttt{DAIR3.sh} --- Project launcher (Linux/WSL)}]
#!/bin/bash
# DAIR3.sh
# Usage: source ~/DAIR3/DAIR3.sh
# Launches the DAIR3 project environment.

source ~/penv/ActivateEnv.sh DAIR3 "$HOME/DAIR3"
\end{lstlisting}

\subsection{Making a Script Executable}
\label{sec:chmod}

On Linux, creating a file does not automatically make it executable. You
must grant execute permission using the \texttt{chmod} command:

\begin{lstlisting}[style=linuxStyle]
chmod +x ~/penv/ActivateEnv.sh
chmod +x ~/DAIR3/DAIR3.sh
\end{lstlisting}

\texttt{chmod} stands for \emph{change mode}. The \texttt{+x} flag adds
execute (\texttt{x}) permission. Linux permissions are organized around
three roles---owner, group, and others---but for personal scripts,
\texttt{chmod +x} is almost always all you need.

\subsection{Adding Your Scripts to PATH}

On Windows, you added \texttt{C:\textbackslash penv} to the system
\texttt{PATH} so that \texttt{ActivateEnv.bat} could be found from
anywhere. On Linux, you do the same by editing your shell configuration
file.

\begin{lstlisting}[style=linuxStyle, caption={Adding \textasciitilde/penv to PATH in \textasciitilde/.bashrc}]
# Open your bashrc file in a text editor
nano ~/.bashrc

# Add this line at the bottom:
export PATH="$HOME/penv:$PATH"

# Save and reload the configuration
source ~/.bashrc
\end{lstlisting}

After this, the shell can find \texttt{ActivateEnv.sh} from any directory.

% ============================================================
\section{The Critical Difference: Executing vs.\ Sourcing}
\label{sec:source-vs-execute}
% ============================================================

This section explains one of the most important---and most
misunderstood---concepts in Linux shell scripting, especially for
developers coming from Windows.

\subsection{Child Shells}

Every time you run a script with \texttt{./script.sh} or just
\texttt{script.sh}, Linux launches a \textbf{child shell}: a new,
separate shell process that inherits a copy of your environment, runs
the script, and then \emph{terminates}. Any changes the script makes to
the environment---activating a venv, changing directory, setting a
variable---exist only in that child shell. When the child terminates,
those changes disappear. Your original terminal is completely unaffected.

This is why running \texttt{./ActivateEnv.sh} would seem to do
nothing useful: the venv gets activated inside the child, then the child
exits, and your terminal is still pointing at the global Python.

\subsection{Sourcing: Running in the Current Shell}

The \texttt{source} command (or its shorthand, a single dot \texttt{.})
runs a script \emph{in the current shell}---no child is created. Every
command executes as if you had typed it directly at the prompt. This means:

\begin{itemize}
  \item \texttt{source ~/penv/DAIR3/bin/activate} activates the venv
        in \emph{your} terminal.
  \item \texttt{cd} changes \emph{your} working directory.
  \item Variables set by the script persist in \emph{your} session.
\end{itemize}

\begin{lstlisting}[style=linuxStyle]
# These two are equivalent:
source ActivateEnv.sh DAIR3 ~/DAIR3
. ActivateEnv.sh DAIR3 ~/DAIR3

# This does NOT work for venv activation:
./ActivateEnv.sh DAIR3 ~/DAIR3    # child shell; changes are lost
\end{lstlisting}

\subsection{The Windows Parallel}

Windows does not have this distinction in quite the same way. The
\texttt{call} command in a batch file runs another \texttt{.bat} in the
same CMD session and returns, which is why the Windows version works
as expected. Linux's default behavior of spawning a child process is
actually the \emph{safer} design---it prevents scripts from accidentally
modifying your environment---but it requires you to opt in to the
shared-session behavior with \texttt{source}.

\begin{notebox}
This is also why \texttt{ActivateEnv.sh} uses \texttt{return} instead of
\texttt{exit} for error handling. When a script is sourced, \texttt{exit}
would close your entire terminal. \texttt{return} exits only the script,
leaving your session intact.
\end{notebox}

% ============================================================
\section{Level 2 --- Shell Functions: The Elegant Linux Approach}
\label{sec:level2}
% ============================================================

Sourcing scripts is powerful, but it requires typing \texttt{source}
or \texttt{.} every time. It also means your launcher script must live
somewhere on the \texttt{PATH} and the user must remember the exact
invocation. Linux offers a more elegant solution: \textbf{shell functions}.

\subsection{What Is a Shell Function?}

A shell function is like a script, but it is defined directly inside
your shell configuration file (\texttt{\textasciitilde/.bashrc}).
Because it is part of the shell itself---not a separate file---it
automatically runs in the current session without requiring
\texttt{source}. You call it just by typing its name.

\subsection{Defining the ActivateEnv Function}

Add the following to the bottom of \texttt{\textasciitilde/.bashrc}:

\begin{lstlisting}[style=linuxStyle, basicstyle=\footnotesize\ttfamily, caption={Shell function version of \texttt{ActivateEnv} in \textasciitilde/.bashrc}]
# -------------------------------------------------------
# ActivateEnv: Activate a Python venv and open VS Code
# Usage: ActivateEnv <ENV_NAME> <PROJECT_DIR>
# -------------------------------------------------------
ActivateEnv() {
    local ENV_NAME="$1"
    local BASE_DIR="$2"

    echo "Starting activation of environment '${ENV_NAME}'..."

    # Validate input
    if [ -z "$ENV_NAME" ]; then
        echo "[ERROR] Environment name not specified."
        return 1
    fi

    if [ -z "$BASE_DIR" ]; then
        echo "[ERROR] Base directory not specified."
        return 1
    fi

    if [ ! -d "$BASE_DIR" ]; then
        echo "[ERROR] Directory not found: ${BASE_DIR}"
        return 1
    fi

    local VENV_ACTIVATE="$HOME/penv/${ENV_NAME}/bin/activate"

    if [ ! -f "$VENV_ACTIVATE" ]; then
        echo "[ERROR] Virtual environment not found: ${VENV_ACTIVATE}"
        return 1
    fi

    # Change directory, activate, and open VS Code
    cd "$BASE_DIR" || return 1
    source "$VENV_ACTIVATE"
    echo "Environment '${ENV_NAME}' activated."
    code .
}
\end{lstlisting}

The keyword \texttt{local} restricts variables to the function's scope,
so they do not pollute your shell session.

\subsection{Defining the DAIR3 Function}

Now add a one-line function that calls \texttt{ActivateEnv} for the
DAIR3 project:

\begin{lstlisting}[style=linuxStyle, caption={Project-specific function for DAIR3 in \textasciitilde/.bashrc}]
# -------------------------------------------------------
# DAIR3: Launch the DAIR3 project environment
# Usage: DAIR3
# -------------------------------------------------------
DAIR3() {
    ActivateEnv "DAIR3" "$HOME/DAIR3"
}
\end{lstlisting}

\subsection{Activating the New Functions}

After editing \texttt{\textasciitilde/.bashrc}, reload it:

\begin{lstlisting}[style=linuxStyle]
source ~/.bashrc
\end{lstlisting}

Now, from \emph{any directory}, simply type:

\begin{lstlisting}[style=linuxStyle]
DAIR3
\end{lstlisting}

The shell will change to your DAIR3 project directory, activate the
virtual environment, and open VS Code---all in one word, with no
\texttt{source} prefix required.

\subsection{Why Windows Has No Clean Equivalent}

Windows batch files always run as child processes; there is no
\texttt{source} command in CMD. PowerShell comes closest: you can define
functions in your PowerShell profile (\texttt{\$PROFILE}), and they
behave similarly to Bash functions. However, PowerShell's profile is
disabled by default for security reasons and requires changing the
execution policy. For Windows users who want this one-command experience
natively, configuring PowerShell profiles is the path forward---but it
involves more administrative overhead than editing
\texttt{\textasciitilde/.bashrc}. This is one of the practical reasons
many developers on Windows adopt WSL.

% ============================================================
\section{Putting It All Together}
\label{sec:together}
% ============================================================

Figure~\ref{fig:workflow} summarizes the complete workflow on each platform.

\begin{table}[h]
\centering
\caption{Complete workflow comparison: Windows vs.\ Linux/WSL}
\label{fig:workflow}
\begin{tabular}{@{}p{0.15\textwidth}p{0.38\textwidth}p{0.38\textwidth}@{}}
\toprule
\textbf{Step} & \textbf{Windows} & \textbf{Linux / WSL} \\
\midrule
Create venv
  & \texttt{python -m venv C:\textbackslash penv\textbackslash DAIR3}
  & \texttt{python3 -m venv \textasciitilde/penv/DAIR3} \\[4pt]
Reusable launcher
  & \texttt{ActivateEnv.bat} in \texttt{C:\textbackslash penv}
  & Function \texttt{ActivateEnv()} in \texttt{\textasciitilde/.bashrc} \\[4pt]
Project launcher
  & \texttt{DAIR3.bat} in project folder
  & Function \texttt{DAIR3()} in \texttt{\textasciitilde/.bashrc} \\[4pt]
Make available globally
  & Add \texttt{C:\textbackslash penv} to system PATH
  & Functions in \texttt{.bashrc} are always available \\[4pt]
Launch project
  & Type \texttt{DAIR3} in any CMD window
  & Type \texttt{DAIR3} in any terminal \\[4pt]
Result
  & venv active, VS Code opens in project folder
  & venv active, VS Code opens in project folder (WSL mode) \\
\bottomrule
\end{tabular}
\end{table}

\subsection{The Full WSL Experience}

When you type \texttt{DAIR3} in your WSL terminal:

\begin{enumerate}
  \item Bash looks up \texttt{DAIR3} in its list of known functions.
  \item \texttt{DAIR3()} calls \texttt{ActivateEnv "DAIR3" "\$HOME/DAIR3"}.
  \item \texttt{ActivateEnv()} changes to the project directory.
  \item The virtual environment is activated in your current shell.
  \item \texttt{code .} launches VS Code with the WSL extension,
        connecting your Windows VS Code interface to your Linux backend.
  \item You are now working in a fully configured development environment
        with one word.
\end{enumerate}

% ============================================================
\section{Quick Reference Cheat Sheet}
\label{sec:cheatsheet}
% ============================================================

\begin{table}[h]
\centering
\caption{Shell scripting and virtual environment cheat sheet}
\label{tab:cheatsheet}
\begin{tabular}{@{}p{0.35\textwidth}p{0.28\textwidth}p{0.28\textwidth}@{}}
\toprule
\textbf{Task} & \textbf{Windows (CMD)} & \textbf{Linux / WSL (Bash)} \\
\midrule
Create virtual environment
  & \texttt{python -m venv C:\textbackslash penv\textbackslash NAME}
  & \texttt{python3 -m venv \textasciitilde/penv/NAME} \\[4pt]
Activate environment
  & \texttt{C:\textbackslash penv\textbackslash NAME\textbackslash Scripts\textbackslash activate}
  & \texttt{source \textasciitilde/penv/NAME/bin/activate} \\[4pt]
Deactivate environment
  & \texttt{deactivate}
  & \texttt{deactivate} \\[4pt]
Install a package
  & \texttt{pip install pkg}
  & \texttt{pip install pkg} \\[4pt]
Save package list
  & \texttt{pip freeze > req.txt}
  & \texttt{pip freeze > req.txt} \\[4pt]
Restore from list
  & \texttt{pip install -r req.txt}
  & \texttt{pip install -r req.txt} \\[4pt]
Make script executable
  & (not required; extension handles it)
  & \texttt{chmod +x script.sh} \\[4pt]
Run a script
  & \texttt{script.bat}
  & \texttt{./script.sh} \\[4pt]
Source a script
  & \texttt{call script.bat}
  & \texttt{source script.sh} \\[4pt]
Edit shell config
  & (no equivalent in CMD)
  & \texttt{nano \textasciitilde/.bashrc} \\[4pt]
Reload shell config
  & (restart CMD)
  & \texttt{source \textasciitilde/.bashrc} \\[4pt]
Open VS Code (WSL)
  & \texttt{code .}
  & \texttt{code .} \\[4pt]
Script parameter~1
  & \texttt{\%\textasciitilde1}
  & \texttt{\$1} \\[4pt]
Check if variable empty
  & \texttt{if "\%VAR\%"==""}
  & \texttt{if [ -z "\$VAR" ]} \\[4pt]
Check if directory exists
  & \texttt{if exist "DIR\textbackslash"}
  & \texttt{if [ -d "DIR" ]} \\[4pt]
Exit a script/function
  & \texttt{goto :eof}
  & \texttt{return} (in function/sourced script) \\
\bottomrule
\end{tabular}
\end{table}

% ============================================================
\section{Summary}
\label{sec:summary}
% ============================================================

This chapter introduced the command-line interface as a text-based,
scriptable, and automatable way to interact with your operating system.
We covered:

\begin{itemize}
  \item The fundamental differences between Windows and Linux CLIs,
        including path conventions, executability, and the \texttt{PATH}
        variable.
  \item How WSL brings a full Linux environment to Windows users, and
        how the VS Code WSL extension creates a unified, professional
        development experience.
  \item Python virtual environments: what they are, why they matter, and
        how to create, activate, populate, and deactivate them on both
        platforms.
  \item \textbf{Level~1} scripting: writing \texttt{ActivateEnv.bat} /
        \texttt{ActivateEnv.sh} and \texttt{DAIR3.bat} / \texttt{DAIR3.sh}
        as reusable launchers.
  \item The critical \texttt{source} vs.\ execute distinction on Linux,
        and why virtual environment activation requires the former.
  \item \textbf{Level~2} scripting: shell functions defined in
        \texttt{\textasciitilde/.bashrc} that provide a one-word launch
        command from any directory, with no \texttt{source} prefix
        required.
\end{itemize}

The \texttt{DAIR3} command you built in this chapter is a small but
meaningful example of a broader principle: the CLI lets you encode your
workflow into repeatable, shareable, and improvable scripts. Every
tedious setup step you eliminate today is cognitive load freed up for the
work that actually matters.

